% 使用 TeX 自带中文字体，避免不同电脑上的系统字体名称差异。
\PassOptionsToPackage{fontset=fandol}{ctex}
\documentclass[withoutpreface,bwprint]{cumcmthesis}
%%%%%%%%%%%%%%%%%%%%%%%%%%%%%%%%%%%%%%%%
\usepackage{newtxtext,newtxmath} % 统一正文与数学公式字体，使全文字体风格一致
\usepackage{listings} % 用于插入和排版 Python、MATLAB、C 等程序代码
\usepackage{etoolbox} % LaTeX 命令修改与补丁工具，用于修改已有命令或环境
\usepackage{tikz} % 用于直接在 LaTeX 中绘制流程图、模型结构图和其他矢量图形
\usepackage{xcolor} % 提供文字、表格、代码、TikZ 等的颜色支持
\usepackage{url} % 用于正确排版和换行网页 URL
\usepackage[framemethod=TikZ]{mdframed} % 创建可分页的文本框、定义框等，使用 TikZ 绘制边框
\usepackage{makecell} % 增强表格单元格功能，支持单元格内换行和多行表头

% 附录代码的统一样式：纯白代码区、四边框和左侧行号。
\lstdefinestyle{appendixcode}{%
  frame=none,
  numbers=left,
  numberstyle=\scriptsize\color{gray!75},
  numbersep=7pt,
  xleftmargin=2.2em,
  xrightmargin=0.5em,
  aboveskip=0pt,
  belowskip=0pt,
  backgroundcolor={},
  showstringspaces=false,
  basicstyle=\small\ttfamily,
  keywordstyle=\color{blue!70!black},
  commentstyle=\color{green!45!black},
  stringstyle=\color{orange!75!black},
  breaklines=true,
  breakatwhitespace=false,
  columns=fullflexible,
  keepspaces=true,
  tabsize=4
}

\newmdenv[
  linewidth=0.75pt,
  linecolor=gray!45,
  topline=true,
  bottomline=true,
  leftline=true,
  rightline=true,
  backgroundcolor=white,
  roundcorner=0pt,
  innerleftmargin=0pt,
  innerrightmargin=0pt,
  innertopmargin=0.6pt,
  innerbottommargin=0.6pt,
  skipabove=0.25em,
  skipbelow=0.8em,
  splittopskip=0.6pt,
  splitbottomskip=0.6pt
]{appendixcodebox}

% mcmthesis v6.3 (local class file) hard-codes red team number in \makesheet; patch it to black.
\makeatletter
\patchcmd{\makesheet}{\textcolor{red}{\MCM@opt@tcn}}{\textcolor{black}{\MCM@opt@tcn}}{}{}
\makeatother

\usepackage{multicol} % 支持将页面内容分成多个并行栏
\usepackage{paracol} % 提供更灵活的多栏排版，并允许不同栏独立排版
\usepackage{indentfirst} % 使章节标题后的第一段也自动首行缩进
\usepackage{booktabs} % 用于制作规范的学术三线表
\usepackage{caption} % 控制图片和表格标题的格式、间距等
\usepackage{subcaption} % 用于制作带有 (a)、(b) 等编号的子图和子表
\usepackage{fancyhdr} % 自定义页眉和页脚

% Header rule (the horizontal line under the header) length is controlled by \headwidth in fancyhdr.
% If you changed margins via geometry, \headwidth may not match \textwidth; force them consistent here.
% You can tweak \headwidth (e.g., 0.95\textwidth) and offsets to fine-tune the visible line length.
\setlength{\headwidth}{\textwidth}
\fancyhfoffset[L]{1pt}
\fancyhfoffset[R]{1pt}

\usepackage{algorithm} % 提供算法浮动体
\usepackage{amsmath,amssymb} % 模板类已提供数学环境；这里只补充常用数学符号
\usepackage{algpseudocode} % 提供 algorithmic 环境及 \State、\While 等伪代码命令
\usepackage{newtxmath} % must come after amsXXX
\usepackage{gbt7714} % 按 GB/T 7714-2015 规范排版 BibTeX 参考文献

% Better line breaking for long URLs (especially in bibliography)
\usepackage{xurl} % 改善长网址的自动换行，尤其适用于参考文献中的 URL
\Urlmuskip=0mu plus 1mu

\bibliographystyle{gbt7714-numerical}

\usepackage{graphicx} % 插入 JPG、PNG、PDF 等图片，并支持缩放、旋转等操作
\usepackage{pdfpages} % 将完整的外部 PDF 页面插入当前论文
\usepackage{multirow} % 允许表格中的一个单元格跨越多行
\usepackage{tabularx} % 创建可自动调整列宽并适应页面宽度的表格
\usepackage{float} % 增强图片、表格等浮动体的位置控制，并提供 [H]
\usepackage{array} % 扩展 LaTeX 表格列格式，支持自定义列类型
\usepackage{ragged2e} % 提供更好的左对齐、右对齐和两端对齐控制
\usepackage{enumitem} % 自定义 itemize 和 enumerate 列表的编号、缩进和间距
\usepackage{titlesec} % 自定义 section、subsection 等标题的格式和上下间距
\usepackage{needspace} % 避免标题及其起始内容越过页脚

% 公式交叉引用按实际编号显示，例如“公式 2”或“公式 5”。
\crefformat{equation}{#2公式#1#3}

% 定义自动伸缩宽度且左对齐的 L 表格列类型
\newcolumntype{L}{>{\raggedright\arraybackslash}X}

% 统一增大三线表的行间距；可将 1.12 调大或调小。
\renewcommand{\arraystretch}{1.12}

% 解决 LaTeX 生成 PDF 时的书签/超链接字符串兼容问题
\pdfstringdefDisableCommands{%
  \def\leavevmode@ifvmode{}%
  \def\kern#1{}%
}

\setlength{\textfloatsep}{4pt plus 2pt minus 2pt}
\setlength{\intextsep}{4pt plus 2pt minus 2pt}
\setlength{\floatsep}{6pt plus 2pt minus 2pt}
\captionsetup{skip=2pt}

% Tighten paragraph + list spacing
% Paragraph spacing helper: adjust this ONE line to change paragraph spacing globally
\newcommand{\SetParSkip}[1]{\setlength{\parskip}{#1}}
\SetParSkip{0.35\baselineskip}

% Heading spacing helpers: adjust section/subsection spacing independently
% #1 is the space *after* the heading (heading -> first paragraph)
\newcommand{\SetSectionAfterSkip}[1]{\titlespacing*{\section}{0pt}{0.2\baselineskip}{#1}}
\newcommand{\SetSubsectionAfterSkip}[1]{\titlespacing*{\subsection}{0pt}{0.15\baselineskip}{#1}}

% Defaults: keep heading-to-text spacing consistent with paragraph spacing
\SetSectionAfterSkip{\parskip}
\SetSubsectionAfterSkip{\parskip}

% 标题不能孤立在页底：空间不足时由 LaTeX 自动转到下一页。
% 这只影响分页，不改变正文版心、页边距或页脚位置。
\pretocmd{\section}{\Needspace{5\baselineskip}}{}{}
\pretocmd{\subsection}{\Needspace{2\baselineskip}}{}{}
\pretocmd{\subsubsection}{\Needspace{\baselineskip}}{}{}

\setlist[itemize]{topsep=0pt, itemsep=0pt, parsep=0pt, partopsep=0pt}
\setlist[enumerate]{topsep=0pt, itemsep=0pt, parsep=0pt, partopsep=0pt}

\usepackage{tocloft} % 自定义目录格式，包括目录字体、缩进和各级标题之间的间距
\makeatletter
\setlength{\cftbeforesecskip}{0.30em}      % 一级标题前的间距
\setlength{\cftbeforesubsecskip}{0.28em}   % 二级标题前的间距
\setlength{\cftbeforesubsubsecskip}{0.26em} % 三级标题前的间距
%\renewcommand{\cftsecfont}{\rmfamily}
%\renewcommand{\cftsubsecfont}{\rmfamily}

%%%%%%%%%%%%%%%%%%%%%%%%%%%%%%%%
\makeatletter
\normalsize
\makeatother

\graphicspath{{.}}  % Place your graphic files in the same directory as your main document
\DeclareGraphicsExtensions{.pdf, .jpg, .tif, .png}

% Increase spacing between the title and the Summary heading on the summary sheet.
% Tune the two values below if you want more/less space.
\makeatletter
\def\make@abstract{%
    \par\vspace*{-0.65cm}% space above Summary (i.e., title → Summary)
    \begin{center}
        \textbf{\abstractname}
    \end{center}
    \vspace*{-0.4cm} % space below Summary (Summary → body)
    \@abstract\par
}
\makeatother

\title{全国大学生数学建模竞赛模板}


\begin{document}
\maketitle
% 摘要页不显示页码；正文从摘要结束后重新从第 1 页开始编号。
\thispagestyle{empty}
\begin{abstract}
\begin{enumerate}
    \item  本文通过研究。。。，提出了。。。
    \item  本文首先给出了。。。的函数，同时对。。。进行了详细的说明，另外引入了。。。来判断。。。是否有效，使得表达更加简洁易懂。
\end{enumerate}


\textbf{对于问题一，}
\begin{enumerate}
    \item 该方案首先基于。。。，设计了一个模型，给出了直接计算。。。的公式。
    \item 本问已经给定了。。。的初始参数，代入模型。。。进行。。。即可得到。。。本文首先证明了。。。，之后通过。。。求得了。。。此外，本文通过。。。进行敏感度分析，发现当。。。（精度），能够兼顾求解的准确性和效率。
\end{enumerate}


\textbf{对于问题二，}

\textbf{对于问题三，}

\textbf{对于问题四，}

\begin{enumerate}
    \item 一句话陈述了本文研究的一般问题
    \item 一两句话总结了整体方法、主要结果和结论
    \item 一两句话回到一般背景，解释比以前的工作取得的进展，并概述针对各个问题展开介绍“方法+结果+优势”，例如我们可以说说自己避免了哪些问题，误差值是多少云云
\end{enumerate}

\keywords{问题一方法  \quad 问题一优化方法  }

\end{abstract}
\setcounter{page}{1}
\pagestyle{plain}

\section{问题重述}

\subsection{问题背景}

\textbf{优化类：}  
（相关定理和关键词的解读）\\
本文针对。。。这一情境设计问题，在。。。（背景条件）给定的条件下，需要综合考虑。。。（可变量）对不同情况下。。。（因变量）的优化，以达到。。。这一目标。

\subsection{问题提出}

基于上述背景，本文建立数学模型解决下列问题：

\textbf{问题1：}  

\textbf{问题2：}  

\textbf{问题3：} 

\textbf{问题4：}  

%%%%%%%%%%%%%%%%%%%%%%%%%%%%%%%%%%%%%%%%%%%%%%%%%%%%%%%%%%%%%

\section{问题分析}

\subsection{问题一分析}
对于问题一，

\subsection{问题二分析}	
对于问题二，

\subsection{问题三分析}
对于问题三，

\subsection{问题四分析}
对于问题四，

如\cref{fig:sample-figure}，这是一只猫。
\begin{figure}[H]%用H强制固定浮动体位置
\centering
\includegraphics[width=6cm]{../figure/一只猫}
\caption{一直猫}
\label{fig:sample-figure}
\end{figure}
%%%%%%%%%%%%%%%%%%%%%%%%%%%%%%%%%%%%%%%%%%%%%%%%%%%%%%%%%%%%% 

\section{模型假设}
本文模型假设主要包括：

\begin{itemize}[itemindent=2em]
\item 假设1
\item 假设2
\item 假设3
\item 假设4
\end{itemize}

%%%%%%%%%%%%%%%%%%%%%%%%%%%%%%%%%%%%%%%%%%%%%%%%%%%%%%%%%%%%% 

\section{符号说明}
\begin{table}[H]
\centering
\caption{符号说明}
\label{tab:symbols}
\begin{tabularx}{0.98\textwidth}{@{}>{\centering\arraybackslash}p{0.24\textwidth} >{\raggedright\arraybackslash}X >{\centering\arraybackslash}p{0.14\textwidth}@{}}
    \toprule
    \textbf{符号} & \textbf{含义} & \textbf{单位} \\\midrule
    $I$ & 机器人集合，索引为 $i$ & -- \\
    $J$ & 任务集合，索引为 $j,l$ & -- \\
    $R$ & 充电站集合，索引为 $r$ & -- \\
    $H$ & 共享资源集合，索引为 $h$ & -- \\
    $\Omega$ & 运行情景集合，索引为 $\omega$ & -- \\
    $x_{ij}$ & 机器人 $i$ 是否参与任务 $j$ 的0--1变量 & 0--1 \\
    $u_{ijkl}$ & 机器人 $i$ 执行任务 $j$ 后执行任务 $l$ 的顺序变量 & 0--1 \\
    $s_j,c_j,d_j$ & 任务 $j$ 的开始时间、完成时间和延期量 & min \\
    $e_{it}$ & 机器人 $i$ 在事件 $t$ 后的电量 & kWh \\
    $q_{ir}$ & 机器人 $i$ 在充电站 $r$ 获得的电量 & kWh \\
    $y_B$ & 是否启用操作员B的0--1变量 & 0--1 \\
    $a_{ij}$ & 机器人 $i$ 对任务 $j$ 的静态可行参数 & -- \\
    \makecell[c]{$work_j,release_j,$\\$due_j$} & 任务工作时间、释放时刻和截止时刻 & min \\
    $load_j$ & 任务负载 & kg \\
    $B_i,\underline e_i,e_{i0}$ & 电池容量、最低电量和初始电量 & kWh \\
    $\alpha_i,\rho_i,\rho_r$ & 能耗率、机器人充电速率和站点充电速率 & kWh/min \\
    $W_i$ & 机器人累计作业时间 & min \\
    $P^{max}$ & 系统总充电功率上限，本文设为0.20 & kWh/min \\
    $p_\omega$ & 情景 $\omega$ 的发生概率 & -- \\\bottomrule
\end{tabularx}
\end{table}

%%%%%%%%%%%%%%%%%%%%%%%%%%%%%%%%%%%%%%%%%%%%%%%%%%%%%%%%%%%%% 

\section{模型建立与求解}

\subsection{问题一模型的建立与求解}

\subsubsection{决策变量与静态资格}

\cref{eq1}
\begin{center}

\begin{equation}\label{eq1}
\begin{aligned}
&\min Z = X_1 + X_2 \\
&\text{s.t.} \quad 
\left\{
\begin{aligned}
A x &\leq b \\
A_{\mathrm{eq}} x &= b_{\mathrm{eq}} \\
lb &\leq x \leq ub
\end{aligned}
\right.
\end{aligned}
\end{equation}
\end{center}

\begin{algorithm}[H]
\caption{Optimal Order Quantity Simulation for Newsvendor Problem}
\label{alg:newsvendor}
\begin{algorithmic}[1]
\Require Unit price $a$, unit cost $b$, Poisson parameter $\lambda$, number of simulations $N$
\Ensure Optimal order quantity $u^*$
\State \textbf{Initialization:}
\State $M_{\text{prev}} \gets 0$ \Comment{Expected profit for order quantity $u=0$}
\State $u \gets 1$ \Comment{Start testing from order quantity $u=1$}
\While{True}
    \State $\text{total\_profit} \gets 0$
    \For{$i \gets 1$ \textbf{to} $N$} \Comment{Simulation phase}
        \State Generate random demand $X \sim \text{Poisson}(\lambda)$
        \State $\text{sales} \gets \min(X, u)$
        \State $\text{profit} \gets a \times \text{sales} - b \times u$
        \State $\text{total\_profit} \gets \text{total\_profit} + \text{profit}$
    \EndFor
    \State $M_{\text{current}} \gets \text{total\_profit} / N$ \Comment{Calculate average expected profit}
    \If{$condition$}
        \State $ $ \Comment{Optimal quantity found}
        \State \textbf{break} \Comment{Exit the loop}
    \Else
        \State $M_{\text{prev}} \gets M_{\text{current}}$ \Comment{Update previous profit}
        \State $u \gets u + 1$ \Comment{Test next order quantity}
    \EndIf
\EndWhile
\Statex
\State \textbf{return} $u^*$
\end{algorithmic}
\end{algorithm}

\subsubsection{约束与目标函数}

\subsubsection{求解结果}

\subsection{问题二模型的建立与求解}

首先，我们使用了\cref{capet}。
\begin{equation}\label{capet}
CapE(t)=\frac{Q_{e,max}}{1+e^{\!\left(-k_e\,(t-t_{e0})\right)}}
\end{equation}

\subsubsection{时间窗与任务不重叠}

\subsubsection{协同与能量状态}

\subsubsection{目标函数与正常调度结果}

\subsection{问题三模型的建立与求解}

\subsubsection{前置关系与共享资源}

\subsubsection{四类情景}

\subsubsection{情景结果}

\subsubsection{情景恢复的求解流程}

\subsection{问题四模型的建立与求解}

\subsubsection{分时电价与充电功率模型}

\subsubsection{操作员覆盖与三种策略}

\subsubsection{统一基准设置}

\subsubsection{策略结果}

\subsubsection{最优策略存在性}

\subsubsection{参数检查与结果一致性}

%%%%%%%%%%%%%%%%%%%%%%%%%%%%%%%%%%%%%%%%%%%%%%%%%%%%%%%%%%%%% 

\section{敏感性分析与稳健性检验}

\subsection{扰动设置}

\subsection{结果与解释}

%%%%%%%%%%%%%%%%%%%%%%%%%%%%%%%%%%%%%%%%%%%%%%%%%%%%%%%%%%%%% 

\section{模型评价与推广}


\subsection{模型优点}
\begin{itemize}[itemindent=2em]
\item 优点1
\item 优点2
\item 优点3
\end{itemize}

\subsection{模型缺点}
\begin{itemize}[itemindent=2em]
\item 缺点1
\item 缺点2
\end{itemize}

\subsection{未来推广}
\begin{itemize}[itemindent=2em]
\item 展望1
\item 展望2
\end{itemize}

%%%%%%%%%%%%%%%%%%%%%%%%%%%%%%%%%%%%%%%%%%%%%%%%%%%%%%%%%%%%% 

\section{AI工具使用声明}

本参赛队在竞赛过程中使用了AI工具，主要用于代码调试和论文语言润色，详细使用情况见支撑材料。

%%%%%%%%%%%%%%%%%%%%%%%%%%%%%%%%%%%%%%%%%%%%%%%%%%%%%%%%%%%%% 
%% 参考文献
\nocite{*}
\bibliography{references}
\newpage
%%%%%%%%%%%%%%%%%%%%%%%%%%%%%%%%%%%%%%%%%%%%%%%%%%%%%%%%%%%%%
%% 附录
\begin{appendices}
% 附录中的表格重新从“表 1”开始编号，不承接正文表格编号。
\setcounter{table}{0}
\renewcommand{\thetable}{\arabic{table}}
    
\section{文件列表}

\begin{table}[H]
\centering
\caption{支撑文件目录}
\begin{tabularx}{0.98\textwidth}{@{}>{\centering\arraybackslash}p{0.4\textwidth} >{\centering\arraybackslash}p{0.6\textwidth}@{}}
\toprule
支撑文件名称   & 文件内容 \\
\midrule
P1.py & 问题一程序代码 \\
P2.py & 问题二程序代码 \\
P3.py & 问题三程序代码 \\
P4.py & 问题四程序代码 \\
\bottomrule
\end{tabularx}
\label{tab:支撑文件目录}
\end{table}

\vspace{0.8em}

\begin{table}[H]
\centering
\caption{附录目录}
\begin{tabularx}{0.98\textwidth}{@{}>{\centering\arraybackslash}p{0.4\textwidth} >{\centering\arraybackslash}p{0.6\textwidth}@{}}
\toprule
文件名   & 功能描述 \\
\midrule
P1.py & 问题一程序代码 \\
P2.py & 问题二程序代码 \\
P3.py & 问题三程序代码 \\
P4.py & 问题四程序代码 \\
\bottomrule
\end{tabularx}
\label{tab:附录目录}
\end{table}

\section{代码}

\noindent\normalsize 附录 1：TOPSIS 综合评价程序（完整运行入口）\par


\begin{appendixcodebox}
\lstinputlisting[style=appendixcode,language=matlab]{../代码/topsis.m}
\end{appendixcodebox}

\end{appendices}

\end{document} 
